% SDRTerm — Weather-satellite pass planner + LRPT/APT capture walkthrough
%
% Build with:  pdflatex pipeline.tex   (twice, so ToC + refs resolve)

\documentclass[11pt,a4paper]{article}

\usepackage[margin=2cm]{geometry}
\usepackage{amsmath,amssymb,amsthm}
\usepackage{tikz}
\usepackage{pgfplots}
\pgfplotsset{compat=1.18}
\usetikzlibrary{arrows.meta, positioning, shapes.geometric, calc,
                decorations.pathreplacing, backgrounds, fit, patterns}
\usepackage{listings}
\usepackage{xcolor}
\usepackage{hyperref}
\usepackage{caption}
\usepackage{booktabs}
\usepackage{siunitx}
\usepackage{parskip}
\usepackage{eurosym}      % for \EUR
\newcommand{\EUR}[1]{\officialeuro\,#1}

\hypersetup{
  colorlinks=true,
  linkcolor=blue!70!black,
  urlcolor=blue!60!black,
  citecolor=black,
}

\lstset{
  basicstyle=\ttfamily\small,
  frame=single,
  breaklines=true,
  showstringspaces=false,
  columns=fullflexible,
  keepspaces=true,
  backgroundcolor=\color{gray!5},
}

\title{Weather Satellites at 137\,MHz\\
       \large A Walkthrough of SDRTerm's \texttt{meteor} Plugin\\
       (For Amateurs Who Want Pictures of Clouds from Space)}
\author{SDRTerm project}
\date{\today}

\begin{document}
\maketitle
\tableofcontents
\vspace{2em}

%═══════════════════════════════════════════════════════════════════
\section{The Big Picture in One Paragraph}
%═══════════════════════════════════════════════════════════════════

Six low-orbit weather satellites regularly fly over your head
transmitting live pictures of the Earth on frequencies near 137\,MHz.
With a about \EUR{25} antenna, a filter, and a hobby SDR you can
receive those pictures directly — no internet, no subscription, no
ground station.  This plugin schedules the passes, retunes your SDR
at rise, hands the samples off to \texttt{satdump} (a battle-tested
external decoder that turns 137\,MHz IQ into PNG images), and stops
the recording at set.  What follows is how all of that hangs together.

\begin{center}
\begin{tikzpicture}[node distance=0.4cm,
  block/.style={draw, rounded corners, minimum width=1.9cm, minimum height=0.9cm,
                align=center, font=\scriptsize\ttfamily},
  arrow/.style={-{Latex[length=2mm]}, thick}]

  \node[block] (celes)              {CelesTrak\\TLE download};
  \node[block, right=of celes] (pred){pyorbital\\pass predict};
  \node[block, right=of pred] (tune) {SDR retune\\on rise};
  \node[block, right=of tune] (tcp)  {rtl-tcp\\passive server};
  \node[block, right=of tcp]  (sd)   {satdump\\decode};
  \node[block, right=of sd]   (png)  {PNG\\image};

  \draw[arrow] (celes) -- (pred);
  \draw[arrow] (pred)  -- (tune);
  \draw[arrow] (tune)  -- (tcp);
  \draw[arrow] (tcp)   -- (sd);
  \draw[arrow] (sd)    -- (png);
\end{tikzpicture}
\end{center}

The plugin owns the four blocks on the left; \texttt{satdump} owns the
DSP and image-reassembly work on the right.  Splitting the work this
way keeps SDRTerm's codebase small (satdump is $\sim$200 k lines of
C++ implementing every mainstream satellite pipeline) and lets us
benefit from bug fixes the wider community makes.

%═══════════════════════════════════════════════════════════════════
\section{The Satellites}
%═══════════════════════════════════════════════════════════════════

Six spacecraft, two families, all in low-Earth polar orbit.  All
downlink near 137\,MHz so \emph{one antenna serves all six}.

\begin{center}
\begin{tabular}{@{}lllll@{}}
\toprule
\textbf{Satellite} & \textbf{Frequency} & \textbf{Mode} & \textbf{Bandwidth} & \textbf{Status (2026)} \\
\midrule
NOAA 15         & \SI{137.620}{MHz} & APT   & \SI{34}{kHz}  & Intermittent since 2019 \\
NOAA 18         & \SI{137.9125}{MHz}& APT   & \SI{34}{kHz}  & Active \\
NOAA 19         & \SI{137.100}{MHz} & APT   & \SI{34}{kHz}  & Active (shares w/ Meteor-M2 3) \\
METEOR-M2 2     & \SI{137.900}{MHz} & LRPT  & \SI{150}{kHz} & Degraded, sporadic \\
METEOR-M2 3     & \SI{137.100}{MHz} & LRPT  & \SI{150}{kHz} & Active \\
METEOR-M2 4     & \SI{137.900}{MHz} & LRPT  & \SI{150}{kHz} & Active \\
\bottomrule
\end{tabular}
\end{center}

The two families broadcast \emph{completely different} kinds of signal:

\begin{itemize}
  \item \textbf{NOAA POES} → \textbf{APT} (Automatic Picture
        Transmission) — an \emph{analog} FM signal carrying an AM
        subcarrier at \SI{2400}{Hz}.  Format hasn't changed since
        1978.  The received picture is a slow-scan raster, two lines
        per second, 909 pixels per line, low resolution but robust —
        even a \SI{10}{dB} SNR link produces a recognisable image.
  \item \textbf{Meteor-M2} → \textbf{LRPT} (Low-Rate Picture
        Transmission) — a modern \emph{digital} QPSK link at
        \SI{72}{ksym/s} with convolutional + Reed-Solomon FEC.
        Pictures are \SI{1568}{pixels} wide across three channels
        (visible + two IR), much higher quality than APT — but the
        signal falls off a cliff when SNR drops below the FEC
        threshold ($\sim$\SI{6}{dB}).
\end{itemize}

\begin{center}
\begin{tikzpicture}[scale=0.85]
  \begin{axis}[width=13cm, height=4cm,
    ymin=-90, ymax=5,
    xlabel={frequency [\si{kHz}]}, ylabel={dBFS},
    axis y line=left, axis x line=bottom,
    xmin=-100, xmax=100, xtick={-100,-50,0,50,100}]
    % APT: narrow FM-shaped hump ~34 kHz wide, centred at 0
    \addplot[very thick, blue, domain=-25:25, samples=100]
        {-15 - 20 * abs(x/25)};
    % LRPT: wider flat-top QPSK spectrum ~150 kHz wide
    \addplot[very thick, red, domain=-75:75, samples=100]
        {-25 - 5*(x/75)^4};
    \node[blue, font=\scriptsize] at (axis cs:-30, -25) {APT (NOAA)};
    \node[red,  font=\scriptsize] at (axis cs: 65, -35) {LRPT (Meteor)};
  \end{axis}
\end{tikzpicture}
\end{center}

%═══════════════════════════════════════════════════════════════════
\section{Orbital Mechanics: Why You Only Get 10 Minutes}
%═══════════════════════════════════════════════════════════════════

\subsection{The orbit}

All six satellites fly \emph{sun-synchronous polar} orbits at
altitudes of \(\sim\)\SI{820}{km}.  Sun-synchronous means the plane
of the orbit rotates \(\sim\)\SI{1}{deg/day} in the same direction
the Earth orbits the Sun, so the satellite crosses the equator at
the same local solar time on every pass.  Polar means the
inclination is close to 98°, so the satellite covers essentially the
whole planet in a day.

Orbital period at 820\,km:
\begin{equation}
T \;=\; 2\pi \sqrt{\frac{(R_\oplus + h)^3}{G M_\oplus}}
\;\approx\; \SI{101}{min}.
\end{equation}

At that altitude the satellite's ground speed is
\(\sim\)\SI{7.4}{km/s}.  From horizon to horizon at a good pass
(overhead) takes about \SI{15}{min} of geometry; the \emph{usable}
window where elevation stays above \(\sim\)10° is closer to
\SI{8}{min}--\SI{12}{min}.

\begin{center}
\begin{tikzpicture}[scale=0.85]
  \begin{axis}[width=13cm, height=5cm,
    ymin=0, ymax=95,
    xlabel={time from AOS [minutes]}, ylabel={elevation [degrees]},
    axis y line=left, axis x line=bottom,
    xmin=0, xmax=15,
    xtick={0,3,6,9,12,15}, ytick={0,15,30,60,90}]
    % Symmetric hump reaching ~80° at ~7.5 min
    \addplot[thick, blue, domain=0:15, samples=200]
        {80 * sin(deg(x/15 * pi))};
    \draw[dashed, red] (axis cs:0,15) -- (axis cs:15,15);
    \node[red, font=\scriptsize] at (axis cs:1.5,18) {min elevation (default 15°)};
    \node[font=\scriptsize] at (axis cs:1.5,4) {AOS};
    \node[font=\scriptsize] at (axis cs:7.5,90) {TCA};
    \node[font=\scriptsize] at (axis cs:13.5,4) {LOS};
  \end{axis}
\end{tikzpicture}
\end{center}

\textbf{Vocabulary}:
\begin{itemize}
  \item \textbf{AOS} (Acquisition of Signal) — satellite rises above
        the horizon at your site.
  \item \textbf{TCA} (Time of Closest Approach) — highest elevation
        of the pass; also the closest range and the strongest signal.
  \item \textbf{LOS} (Loss of Signal) — satellite sets below the horizon.
  \item \textbf{Elevation} — angle above the local horizon, 0° at
        horizon, 90° straight up.
  \item \textbf{Azimuth} — compass bearing where the satellite is,
        0° = north, 90° = east, etc.
\end{itemize}

Only \emph{high-elevation} passes give clean images.  The plugin
defaults to \texttt{min\_elevation\_deg = 15.0} so you don't waste
storage on horizon-scraping passes that will decode to noise.

\subsection{Where the plugin gets the numbers}

For each satellite we need to solve: \emph{given today's date and
your location, when does the satellite rise, when does it set, and
how high does it get in between?}  Two ingredients:

\begin{enumerate}
  \item \textbf{TLE (Two-Line Element set)} — a compact ASCII
        description of the satellite's orbit at a reference epoch.
        Every satellite of interest is published daily by NORAD via
        CelesTrak.  We fetch the ``weather'' group bundle (about
        \SI{14}{kB}, $\sim$200 satellites) and cache it locally for
        12\,h.
  \item \textbf{SGP4} — an orbital propagator that turns a TLE at
        epoch \(t_0\) into position \(\vec r(t)\) and velocity
        \(\vec v(t)\) at any nearby time \(t\).  We use
        \texttt{pyorbital.Orbital}, which wraps SGP4 in a friendly
        Python API.
\end{enumerate}

An example TLE for NOAA 19:
\begin{lstlisting}
NOAA 19
1 33591U 09005A   26224.53275463  .00000242  00000-0  15012-3 0  9995
2 33591  99.0193 264.7583 0013821 265.4837  94.4569 14.13141235894382
\end{lstlisting}

That's the whole orbit: inclination 99.02°, mean motion \SI{14.13}{rev/day},
eccentricity 0.00138, etc.  From those numbers plus the reference
time in the first line, SGP4 gives us position and elevation at any
requested moment.  The plugin evaluates elevation every
\SI{30}{seconds} for the next 24\,h and picks out contiguous
above-threshold runs — those are our passes.

\begin{lstlisting}[language=Python]
# plugins/meteor/passes.py — the entire prediction is a few lines
o = Orbital(sat_name, tle_file=tle_path)
lat, lon, alt = self._location_lat, self._location_lon, self._location_alt_km
elevs = [o.get_observer_look(t, lon, lat, alt)[1] for t in times]
# Then find contiguous elevs[i] >= min_elevation_deg → passes
\end{lstlisting}

%═══════════════════════════════════════════════════════════════════
\section{The Antenna and the RF Front End}
%═══════════════════════════════════════════════════════════════════

\subsection{Why 137\,MHz needs a special antenna}

A wire whip receives \emph{linearly-polarised} signals.  Both APT
and LRPT are \emph{right-hand circularly-polarised} (RHCP) — the
electric field vector rotates as it propagates, tracing out a
right-handed helix.  A linear antenna captures roughly half the
energy from a circular wave (a \SI{3}{dB} penalty), and worse, it
suffers deep nulls when the satellite is at certain azimuths — which
happens exactly when the sat is passing overhead and \emph{should}
be strongest.

The standard receive antenna for 137\,MHz is a \textbf{Quadrifilar
Helix} (QFH): four helical wires wound around a common axis, fed in
quadrature.  Overall it's a narrow \SI{130}{cm}-tall cage that has:

\begin{itemize}
  \item Nearly-hemispherical gain pattern (peaks upward, drops
        gracefully toward the horizon)
  \item Native RHCP polarisation — matched to the satellites
  \item No nulls in the upper hemisphere (unlike a Yagi or a
        turnstile)
\end{itemize}

You can build a QFH from copper pipe for about about \EUR{25}, or
buy a pre-cut one for \EUR{60} to \EUR{100}.

\begin{center}
\begin{tikzpicture}[scale=0.9]
  % QFH cartoon — two crossed helices
  \draw[thick, blue] plot [domain=0:6.28, samples=100]
      ({0.6*cos(deg(\x))}, {0.4*sin(deg(\x))+\x/2});
  \draw[thick, red, dashed] plot [domain=0:6.28, samples=100]
      ({0.6*cos(deg(\x)+90)}, {0.4*sin(deg(\x)+90)+\x/2});
  \draw[very thick] (0,0) circle (0.6);
  \draw[very thick] (0,3.14) circle (0.6);
  \node[font=\scriptsize, right] at (0.8,3.5) {top loop};
  \node[font=\scriptsize, right] at (0.8,-0.1) {feed / coax};
  \node[font=\scriptsize] at (4,1.5) {\begin{tabular}{l}
     Two crossed helical loops,\\
     $90^\circ$ phase difference\\
     $\Rightarrow$ RHCP, hemispherical\\
     gain pattern.
  \end{tabular}};
\end{tikzpicture}
\end{center}

\subsection{The filter and preamplifier}

The 137\,MHz band sits between two very strong crowded neighbours:

\begin{itemize}
  \item \textbf{Below}: broadcast FM (88--108\,MHz) — 100\,kW ERP
        stations that will drown your SDR's front end in intermod if
        left unfiltered.
  \item \textbf{Above}: airband voice (118--137\,MHz) — dozens of
        AM channels; individually weaker than FM but numerous.
\end{itemize}

A bandpass filter centred at 137\,MHz with a \SI{4}{MHz} passband
kills both crowds.  The \textbf{SAWbird+NOAA} is the standard
combo: SAW filter (\SI{40}{dB} out-of-band rejection) plus a
low-noise amplifier (\SI{20}{dB} gain, \SI{1.5}{dB} noise figure) on
the same tiny board, powered by bias-tee over the coax.

\begin{center}
\begin{tikzpicture}[node distance=0.5cm,
  block/.style={draw, rounded corners, minimum width=1.8cm, minimum height=0.8cm,
                align=center, font=\scriptsize\ttfamily},
  arrow/.style={-{Latex[length=2mm]}, thick}]

  \node[block] (ant)              {QFH\\antenna};
  \node[block, right=of ant]  (saw){SAW BPF\\135--139\,MHz};
  \node[block, right=of saw]  (lna){LNA\\+20\,dB};
  \node[block, right=of lna]  (sdr){HackRF /\\RTL-SDR};

  \draw[arrow] (ant) -- (saw);
  \draw[arrow] (saw) -- (lna);
  \draw[arrow] (lna) -- (sdr);
  \draw[dashed, thick] ($(saw.south west)+(-0.05,-0.05)$) rectangle ($(lna.north east)+(0.05,0.3)$);
  \node[font=\tiny, gray] at ($(saw.north east)+(-0.5,0.5)$) {SAWbird+NOAA};
\end{tikzpicture}
\end{center}

%═══════════════════════════════════════════════════════════════════
\section{APT: The Analog Picture Format}
%═══════════════════════════════════════════════════════════════════

APT has been in operation since 1978 and hasn't changed a bit —
that's why NOAA 15's 1998-vintage transmitter still works.  The
format is beautifully simple.

\subsection{Wire-level modulation}

\begin{enumerate}
  \item \textbf{Video baseband}: a \SI{2400}{Hz} carrier is
        \emph{amplitude-modulated} by the image brightness (dark = low
        AM level, bright = high).
  \item \textbf{FM main link}: that AM subcarrier (plus the
        \SI{2400}{Hz} tone-line telemetry pulses) is
        \emph{frequency-modulated} onto the RF carrier at
        \SI{137}{MHz} with a peak deviation of \SI{17}{kHz}.
        Total occupied bandwidth: about \SI{34}{kHz}.
\end{enumerate}

The decoder therefore has to un-do modulation twice:

\begin{center}
\begin{tikzpicture}[node distance=0.35cm,
  block/.style={draw, rounded corners, minimum width=1.9cm, minimum height=0.7cm,
                align=center, font=\scriptsize\ttfamily},
  arrow/.style={-{Latex[length=2mm]}, thick}]

  \node[block] (iq)              {IQ @\\137.100\,MHz};
  \node[block, right=of iq]  (fm){FM\\demod};
  \node[block, right=of fm]  (am){AM\\envelope};
  \node[block, right=of am]  (sy){sync\\word};
  \node[block, right=of sy]  (ln){line\\assemble};
  \node[block, right=of ln]  (im){grayscale\\raster};

  \draw[arrow] (iq) -- (fm);
  \draw[arrow] (fm) -- (am);
  \draw[arrow] (am) -- (sy);
  \draw[arrow] (sy) -- (ln);
  \draw[arrow] (ln) -- (im);
\end{tikzpicture}
\end{center}

\subsection{Frame structure}

Each APT line is \emph{two sub-lines} stitched together — one from
the visible channel, one from an infrared channel.  Two lines
per second means the receiver produces roughly \SI{60}{lines/minute}.
A whole 10-minute pass gives you $\sim$600 lines, or a $\sim$909 $\times$
600-pixel image (before any oversampling to sharpen it).

\begin{center}
\begin{tikzpicture}[scale=0.75]
  \draw[thick] (0,0) rectangle (16,1);
  \draw[fill=blue!20]  (0,0)   rectangle (1.2,1);   \node[font=\scriptsize] at (0.6,0.5) {SyncA};
  \draw[fill=yellow!30](1.2,0) rectangle (2.6,1);   \node[font=\scriptsize] at (1.9,0.5) {space A};
  \draw[fill=gray!30]  (2.6,0) rectangle (8,1);     \node[font=\scriptsize] at (5.3,0.5) {Video channel A (909 px)};
  \draw[fill=orange!30](8,0)   rectangle (8.4,1);   \node[font=\scriptsize] at (8.2,0.5) {t};
  \draw[fill=blue!20]  (8.4,0) rectangle (9.6,1);   \node[font=\scriptsize] at (9,0.5)   {SyncB};
  \draw[fill=yellow!30](9.6,0) rectangle (11,1);    \node[font=\scriptsize] at (10.3,0.5){space B};
  \draw[fill=gray!30]  (11,0)  rectangle (16,1);    \node[font=\scriptsize] at (13.5,0.5){Video channel B (909 px)};
  \node[font=\scriptsize] at (8,1.4) {one APT line = 0.5\,s = 2080 words @ 4160 words/s};
\end{tikzpicture}
\end{center}

\begin{itemize}
  \item \textbf{Sync words} (7 alternating cycles at \SI{1040}{Hz})
        mark the start of each sub-line — the decoder locks onto these
        with a correlator.
  \item \textbf{Space} pixels are calibration references — cold-target
        levels the receiver uses to normalise brightness across the pass.
  \item \textbf{Telemetry} (the tiny middle segment ``t'') encodes
        which spectral band each channel is using (visible, near-IR,
        thermal IR, water vapour) — different combinations per season.
\end{itemize}

satdump's \texttt{noaa\_apt} pipeline does all of this and outputs
one PNG per channel, plus a false-colour composite.

%═══════════════════════════════════════════════════════════════════
\section{LRPT: The Digital Picture Format}
%═══════════════════════════════════════════════════════════════════

Meteor-M2 uses a much more modern transmission — same era as GPS L2C
or Iridium — but the signal shape and error correction are more
complex.

\subsection{Wire-level modulation}

\begin{itemize}
  \item \textbf{Modulation}: QPSK (Quadrature Phase-Shift Keying) at
        \SI{72}{ksym/s}, so \SI{144}{kbit/s} raw before FEC.
        Occupied bandwidth (root-raised-cosine, roll-off 0.6):
        \(\sim\)\SI{150}{kHz}.
  \item \textbf{Constellation}: 4 points on the unit circle at
        \(\pm 45^\circ, \pm 135^\circ\).  Each symbol carries 2 bits.
  \item \textbf{FEC layer 1}: rate-1/2 convolutional code (constraint
        length 7, generators 171 / 133 octal) → decoded by a Viterbi
        algorithm.
  \item \textbf{FEC layer 2}: Reed-Solomon (255, 223) per packet →
        corrects any remaining byte errors.
  \item \textbf{Transport}: CCSDS (Consultative Committee for Space
        Data Systems) packets, \SI{1024}{bytes} each, addressed to
        virtual channels for different image streams.
\end{itemize}

\begin{center}
\begin{tikzpicture}[scale=1.0]
  \draw[gray!50, ->] (-2,0) -- (2,0) node[right, font=\scriptsize] {I};
  \draw[gray!50, ->] (0,-2) -- (0,2) node[above, font=\scriptsize] {Q};
  \foreach \x/\y/\lbl in {1.2/1.2/{00}, 1.2/-1.2/{10}, -1.2/1.2/{01}, -1.2/-1.2/{11}} {
    \filldraw[fill=blue!30, draw=blue!60!black] (\x,\y) circle (0.2);
    \node[font=\tiny] at (\x*1.35, \y*1.35) {\lbl};
  }
  \draw[dashed] (0,0) circle (1.7);
  \node[font=\scriptsize] at (3.5,1.5) {\begin{tabular}{l}
    QPSK: one of 4 phases per\\
    symbol, each encoding 2 bits.\\
    Baud 72\,000, so \\
    144\,000 raw bits/s.
  \end{tabular}};
\end{tikzpicture}
\end{center}

\subsection{The decoding pipeline (satdump owns this)}

\begin{center}
\begin{tikzpicture}[node distance=0.3cm,
  block/.style={draw, rounded corners, minimum width=1.9cm, minimum height=0.7cm,
                align=center, font=\scriptsize\ttfamily},
  arrow/.style={-{Latex[length=2mm]}, thick}]

  \node[block] (iq)                 {IQ @\\150\,kHz};
  \node[block, right=of iq]   (rrc) {RRC\\match filter};
  \node[block, right=of rrc]  (pll) {Costas\\PLL};
  \node[block, right=of pll]  (sym) {symbol\\timing};
  \node[block, right=of sym]  (vit) {Viterbi\\decode};
  \node[block, right=of vit]  (rs)  {Reed-\\Solomon};
  \node[block, right=of rs]   (ccs) {CCSDS\\depacket};
  \node[block, right=of ccs]  (msu) {MSU-MR\\reassemble};

  \draw[arrow] (iq)  -- (rrc);
  \draw[arrow] (rrc) -- (pll);
  \draw[arrow] (pll) -- (sym);
  \draw[arrow] (sym) -- (vit);
  \draw[arrow] (vit) -- (rs);
  \draw[arrow] (rs)  -- (ccs);
  \draw[arrow] (ccs) -- (msu);
\end{tikzpicture}
\end{center}

\begin{enumerate}
  \item \textbf{RRC match filter} — the transmitter shapes each
        symbol with a root-raised-cosine pulse; the receiver applies
        the matching filter to maximise SNR at the symbol instants.
  \item \textbf{Costas PLL} — locks a local oscillator to the
        received carrier phase; without this, the constellation
        rotates slowly and the decoder decodes garbage.
  \item \textbf{Symbol timing recovery} — finds the exact instants
        (relative to sample clock) when to sample I and Q.  Uses a
        Gardner or M\&M timing-error detector.
  \item \textbf{Viterbi decoder} — un-does the rate-1/2 convolutional
        code by finding the most-likely input sequence through a
        trellis of $2^6 = 64$ states.  Recovers 1 bit for every 2
        received bits.
  \item \textbf{Reed-Solomon decoder} — corrects up to 16 byte
        errors per 255-byte block.  Meteor uses depth-4 interleaving
        so a single burst error is spread across 4 RS blocks.
  \item \textbf{CCSDS depacket} — packet framing, sync detection,
        virtual-channel routing.
  \item \textbf{MSU-MR reassemble} — the six virtual channels contain
        the six spectral bands from the MSU-MR camera; the decoder
        maps packets back into a 2D raster and outputs one PNG per
        channel plus false-colour composites.
\end{enumerate}

%═══════════════════════════════════════════════════════════════════
\section{How the Handoff to satdump Works}
%═══════════════════════════════════════════════════════════════════

SDRTerm does not itself do any of the QPSK/Viterbi/Reed-Solomon work.
Instead, on a predicted pass rise it spawns \texttt{satdump} as a
subprocess and lets it pull the samples over TCP.

\subsection{The rtl-tcp bridge}

SDRTerm ships with an \texttt{rtl-tcp-passive} plugin that mirrors
its captured IQ over the standard \texttt{rtl\_tcp} protocol on
localhost:1234.  \texttt{satdump} was written to talk exactly this
protocol as one of its many source options — so from satdump's point
of view SDRTerm looks like a headless SDR server.

\begin{center}
\begin{tikzpicture}[node distance=0.5cm,
  block/.style={draw, rounded corners, minimum width=2.0cm, minimum height=0.9cm,
                align=center, font=\scriptsize\ttfamily},
  arrow/.style={-{Latex[length=2mm]}, thick},
  proc/.style={fill=gray!5}]

  \node[block, proc] (hw)              {HackRF /\\RTL-SDR};
  \node[block, proc, right=of hw]  (st){SDRTerm\\core + rtl-tcp\\plugin};
  \node[block, proc, right=of st]  (sd){satdump\\subprocess};

  \draw[arrow] (hw) -- node[above, font=\tiny] {USB IQ} (st);
  \draw[arrow] (st) -- node[above, font=\tiny] {TCP :1234} (sd);
  \draw[arrow, dashed] (sd.north) to[bend right=30] node[above, font=\tiny]
       {spawn / signal (fork+exec)} (st.north);
\end{tikzpicture}
\end{center}

Concretely, this is the command line the meteor plugin invokes when
a Meteor-M2 pass begins:

\begin{lstlisting}
satdump live meteor_m2-x_lrpt <output_dir> \
       --source     rtltcp        \
       --frequency  137100000     \
       --samplerate <state.bw_hz> \
       --tcp_address 127.0.0.1    \
       --tcp_port    1234         \
       --timeout    <pass_seconds + 30>
\end{lstlisting}

The \texttt{--timeout} is set to the pass's remaining seconds plus a
30-second slack, so if SDRTerm crashes without cleaning up satdump
gracefully, satdump still exits by itself when the pass would have
ended.

\subsection{Why not do it all in Python?}

Fair question.  The honest answer:

\begin{itemize}
  \item \textbf{Correctness.}  satdump has years of tuning on the
        Meteor-M2 specifics: Doppler compensation, differential QPSK
        variants, degraded-signal recovery, changing packet formats.
        Recreating that in Python would take months and would be worse.
  \item \textbf{Speed.}  Viterbi on 144\,kbit/s at $2^{6}$ trellis
        states + RS-255 on the output is well within Python's reach
        but not with the elegant code style we want in the plugin.
        Delegating keeps SDRTerm's Python footprint focused on
        orchestration.
  \item \textbf{Community.}  Fixes upstream (support for METEOR-M2 4's
        slightly different frame format, for example) reach us for
        free when the user updates satdump.
\end{itemize}

%═══════════════════════════════════════════════════════════════════
\section{The Auto-Tune and Auto-Capture State Machine}
%═══════════════════════════════════════════════════════════════════

Every call to \texttt{process()} (which fires many times per second
in the SDRTerm main loop) walks through this decision tree:

\begin{center}
\begin{tikzpicture}[node distance=0.6cm,
  block/.style={draw, rounded corners, minimum width=3.2cm, minimum height=0.7cm,
                align=center, font=\scriptsize\ttfamily},
  dec/.style={draw, diamond, aspect=2, minimum width=3.2cm, minimum height=0.9cm,
              align=center, font=\scriptsize\ttfamily},
  arrow/.style={-{Latex[length=2mm]}, thick}]

  \node[dec] (any)                     {any pass now?};
  \node[block, right=1cm of any] (kill){if satdump running:\\SIGTERM it, reap};
  \node[dec, below=of any]  (tuned)    {already tuned to sat?};
  \node[block, right=1cm of tuned] (rt){retune SDR to sat freq};
  \node[dec, below=of tuned] (spawned) {satdump already spawned?};
  \node[block, right=1cm of spawned] (sp){spawn satdump};
  \node[block, below=of spawned] (poll){poll satdump.poll();\\reap if exited};

  \draw[arrow] (any) -- node[above, font=\tiny] {no} (kill);
  \draw[arrow] (any) -- node[right, font=\tiny] {yes} (tuned);
  \draw[arrow] (tuned) -- node[above, font=\tiny] {no} (rt);
  \draw[arrow] (tuned) -- node[right, font=\tiny] {yes} (spawned);
  \draw[arrow] (spawned) -- node[above, font=\tiny] {no} (sp);
  \draw[arrow] (spawned) -- node[right, font=\tiny] {yes} (poll);
\end{tikzpicture}
\end{center}

Two invariants keep this from spawning satdump repeatedly or
retuning on every frame:

\begin{itemize}
  \item \texttt{self.\_tuned\_sat / \_tuned\_freq\_hz} — set after a
        successful retune.  Comparing against the current pass's sat
        name short-circuits the retune on subsequent calls.
  \item \texttt{self.\_capture\_sat} — set on \emph{first attempt} to
        spawn satdump, whether or not the spawn succeeded.  Prevents
        us from re-attempting every call when satdump is missing from
        \texttt{\$PATH}.
\end{itemize}

On pass fall, \texttt{find\_current\_pass()} returns \texttt{None}
and the plugin walks the left branch: \texttt{SIGTERM} satdump, wait
up to 10\,s for it to flush its file writes, release the handles.

%═══════════════════════════════════════════════════════════════════
\section{Configuration}
%═══════════════════════════════════════════════════════════════════

Everything is in the preset under
\texttt{plugin\_states.meteor}:

\begin{lstlisting}
"meteor": {
  "location_lat":      49.740693,
  "location_lon":      6.660242,
  "location_alt_km":   0.245,
  "min_elevation_deg": 15.0,
  "horizon_hours":     24.0,
  "auto_tune":         true,
  "auto_capture":      true,
  "rtltcp_host":       "127.0.0.1",
  "rtltcp_port":       1234
}
\end{lstlisting}

Location goes in as decimal degrees (positive = N/E), altitude in
kilometres above mean sea level.  For 137\,MHz reception 100\,m of
altitude error changes reception by nothing meaningful; you can guess
if you don't know it exactly.

\subsection{What happens if you're offline}

TLE fetch fails silently (falls back to the cached copy on disk if
one exists), so an existing installation keeps predicting passes
even without internet.  Fresh installations without any cache will
fail to predict — the first thing the plugin does when it can't
predict passes is tell you why in the web view.

%═══════════════════════════════════════════════════════════════════
\section{Output: Where the Pictures End Up}
%═══════════════════════════════════════════════════════════════════

satdump writes into
\texttt{plugins/meteor/web/captures/<ISO-timestamp>\_<sat>/}
which is under the webserver plugin's static-file tree.  So a
completed pass produces this on disk:

\begin{lstlisting}
plugins/meteor/web/captures/
  2026-08-15T13-42-00Z_METEOR-M2_4/
    METEOR-M2 4_MSU-MR_1.png         ← channel 1 (visible)
    METEOR-M2 4_MSU-MR_2.png         ← channel 2 (visible)
    METEOR-M2 4_MSU-MR_3.png         ← channel 3 (near-IR)
    METEOR-M2 4_MSU-MR_composite.png ← false-colour RGB
    dataset.json                     ← satdump metadata
    satdump.log                      ← stdout+stderr from the run
\end{lstlisting}

The webserver plugin serves these directly at
\texttt{/static/meteor/captures/...} so you can browse them from
any device on your network without opening a file manager on the
Pi.

%═══════════════════════════════════════════════════════════════════
\section{Summary Table}
%═══════════════════════════════════════════════════════════════════

\begin{center}
\begin{tabular}{@{}rll@{}}
\toprule
\textbf{Stage} & \textbf{Owner}                  & \textbf{Purpose} \\
\midrule
1 & CelesTrak                     & TLE feed for weather sats \\
2 & meteor.py + pyorbital        & Predict passes for next 24\,h \\
3 & meteor.py                    & Retune SDR to sat freq on pass rise \\
4 & rtl-tcp-passive plugin       & Serve IQ over TCP :1234 \\
5 & satdump (subprocess)         & Demod + Viterbi + RS + reassemble \\
6 & satdump                      & Write PNGs + metadata to disk \\
7 & webserver plugin             & Serve captures + pass schedule as HTML \\
\bottomrule
\end{tabular}
\end{center}

%═══════════════════════════════════════════════════════════════════
\section{Things We Don't Do (Yet or Ever)}
%═══════════════════════════════════════════════════════════════════

\begin{itemize}
  \item \textbf{Doppler-shift correction on the SDRTerm side.}
        The Doppler shift at 137\,MHz for a satellite in LEO is
        \(\pm\)\SI{3}{kHz} across a pass — enough to matter for LRPT
        (whose entire signal is only 150\,kHz wide).  satdump handles
        it internally using orbital predictions; we don't retune.
  \item \textbf{Antenna rotator control.}  QFH is a fixed
        omnidirectional antenna, so there's nothing to rotate.  For
        higher-gain HRPT reception at L-band you'd need one; that's
        a different plugin entirely.
  \item \textbf{Real-time pass table on the tty.}  Only the web view
        shows the schedule.  Adding a curses tab would be $\sim$50
        lines but nobody's asked yet.
  \item \textbf{Metop / FengYun HRPT decoding.}  Those satellites
        transmit at 1.7\,GHz with much wider bandwidth — completely
        different antenna (\SI{1.2}{m} dish) and much more expensive
        hardware.  Not on the roadmap.
\end{itemize}

%═══════════════════════════════════════════════════════════════════
\section{Bibliography and Further Reading}
%═══════════════════════════════════════════════════════════════════

\begin{itemize}
  \item Vallado, D.\ A., \emph{Fundamentals of Astrodynamics and
        Applications}, 4th ed.  Chapters 8--9 cover SGP4 in detail.
  \item CCSDS 131.0-B-3, \emph{TM Synchronization and Channel
        Coding}.  The convolutional + Reed-Solomon layer used by
        Meteor-M and half of the space-data-transport world.
  \item NOAA POES User Guide, sections on APT format
        (\texttt{https://www.noaasis.noaa.gov/NOAASIS/pubs/}) — the
        original 1978 specification, still authoritative.
  \item \texttt{satdump} project — \texttt{https://www.satdump.org} —
        source of the actual DSP pipelines we invoke.
  \item CelesTrak — \texttt{https://celestrak.org} — for daily TLEs.
  \item pyorbital documentation — SGP4 in a friendly Python wrapper.
\end{itemize}

\end{document}
